\documentclass[12pt]{article}
\usepackage[qed,sharedcounter]{../logicbox}
\newlogicboxtheorem{claim}{Claim}{teal!60!black}{claim}

\title{Logic Box Test}
\author{Test Author}

\begin{document}

\maketitle

\section{Definitions}

\begin{definition}{Group}{def:group}
A \textbf{group} is a set $G$ with a binary operation $\cdot: G \times G \to G$ satisfying:
\begin{enumerate}
    \item Associativity: $(a \cdot b) \cdot c = a \cdot (b \cdot c)$
    \item Identity: $\exists e \in G$ such that $e \cdot a = a \cdot e = a$
    \item Inverse: $\forall a \in G, \exists a^{-1} \in G$ such that $a \cdot a^{-1} = a^{-1} \cdot a = e$
\end{enumerate}
\end{definition}

\section{Theorems}

\begin{axiom}{Line through two points}{ax:line}
For any two points $A$ and $B$, there exists a line passing through both $A$ and $B$.
\end{axiom}

\begin{theorem}{Pythagorean Theorem}{thm:pythagoras}
In a right triangle with legs $a$ and $b$ and hypotenuse $c$, we have $a^2 + b^2 = c^2$.
\end{theorem}

\begin{lemma}{Sum of Angles}{lem:angles}
The sum of interior angles in any triangle is $180^\circ$.
\end{lemma}

\begin{corollary}{Interior Angles of Quadrilateral}{cor:quad}
The sum of interior angles in any quadrilateral is $360^\circ$.
\end{corollary}

\begin{proposition}{Basic Proposition}{prop:even}
If $n$ is an even integer, then $n^2$ is also even.
\end{proposition}

\begin{theorem}{My Theorem}{thm:my}
This is a theorem with optional argument title.
\end{theorem}

\begin{axiom}{Axiom of Choice}{ax:choice}
For any set of non-empty sets, there exists a choice function.
\end{axiom}

\begin{lemma}{Important Lemma}{lem:important}
This is a lemma with an optional argument.
\end{lemma}

\begin{corollary}{Corollary}{cor:my}
This is a corollary with optional title.
\end{corollary}

\begin{proposition}{Proposition}{prop:my}
This is a proposition.
\end{proposition}

\begin{claim}{A Custom Claim}{claim:custom}
This claim uses a user-defined logicbox environment and shares the same section
counter as the theorem-style environments above.
\end{claim}

\section{Proofs}

\begin{proof}
We prove by induction that $\sum_{i=1}^n i = \frac{n(n+1)}{2}$.

\textbf{Base case:} For $n=1$, $\sum_{i=1}^1 i = 1 = \frac{1(1+1)}{2}$.

\textbf{Inductive step:} Assume true for $n$, then
\begin{align*}
\sum_{i=1}^{n+1} i
  &= \sum_{i=1}^n i + (n+1) = \frac{n(n+1)}{2} + (n+1) \\
  &= \frac{n(n+1) + 2(n+1)}{2} = \frac{(n+1)(n+2)}{2}.
\end{align*}
Thus the formula holds for $n+1$, which completes the induction.
\end{proof}

\end{document}
